%% ============================================================
%% Episode 10 - Farsi Podcast Script (Translation)
%% Source: specialized_advisor_report_complete/markdown/10_cross_consistency.md
%% Series: Advisor Progress Report - Deep Dive
%% Compiler: XeLaTeX ONLY
%% Font: B Nazanin
%% CRITICAL: xepersian must always be loaded LAST
%% ============================================================

\documentclass[12pt, a4paper]{article}

%% ============================================================
%% COMPATIBILITY SHIMS
%% ============================================================
\makeatletter
\providecommand\IfClassLoadedT[2]{\@ifclassloaded{#1}{#2}{}}
\providecommand\IfClassLoadedF[2]{\@ifclassloaded{#1}{}{#2}}
\providecommand\IfPackageLoadedT[2]{\@ifpackageloaded{#1}{#2}{}}
\providecommand\IfPackageLoadedF[2]{\@ifpackageloaded{#1}{}{#2}}
\providecommand\IfFileLoadedTF[3]{\@ifpackageloaded{#1}{#2}{#3}}
\providecommand\IfFileLoadedT[2]{\@ifpackageloaded{#1}{#2}{}}
\providecommand\IfFileLoadedF[2]{\@ifpackageloaded{#1}{}{#2}}
\makeatother
\usepackage{expl3}
\ExplSyntaxOn
\cs_if_exist:NF \prop_gput_if_not_in:Nnn
  {
    \cs_new:Npn \prop_gput_if_not_in:Nnn #1#2#3
      { \prop_if_in:NnF #1 {#2} { \prop_gput:Nnn #1 {#2} {#3} } }
  }
\ExplSyntaxOff
%% ============================================================

\usepackage[
  a4paper,
  top=2.5cm, bottom=2.5cm,
  left=2.8cm, right=2.8cm,
  headheight=25pt
]{geometry}

\usepackage{xcolor}
\definecolor{seriesblue}{RGB}{25, 80, 150}
\definecolor{audionote}{RGB}{130, 50, 15}
\definecolor{sectiongray}{RGB}{75, 75, 75}
\definecolor{audiotan}{RGB}{252, 243, 233}

\usepackage{amsmath}
\usepackage{amssymb}
\usepackage{enumitem}
\usepackage{setspace}
\setstretch{1.8}
\usepackage{mdframed}
\usepackage{titlesec}

%% ============================================================
%% xepersian MUST be loaded LAST
%% ============================================================
\usepackage{xepersian}
\settextfont[Scale=1.05]{B Nazanin}
\setlatintextfont{Times New Roman}

\newcommand{\dash}{\lr{---}}

%% ============================================================
%% Page style -- savebox approach
%% ============================================================
\newsavebox{\hdrLeft}
\newsavebox{\hdrRight}
\AtBeginDocument{%
  \savebox{\hdrLeft}{\normalfont\small\color{sectiongray}%
    گزارش پيشرفت مشاور / بررسی عمیق}%
  \savebox{\hdrRight}{\normalfont\small\color{sectiongray}%
    قسمت دهم: سازگاری متقاطع}%
}
\makeatletter
\def\ps@podcast{%
  \def\@oddhead{\usebox{\hdrLeft}\hfill\usebox{\hdrRight}}%
  \let\@evenhead\@oddhead
  \def\@oddfoot{\normalfont\hfil\thepage\hfil\relax}%
  \let\@evenfoot\@oddfoot
  \let\@mkboth\markboth
}
\makeatother
\pagestyle{podcast}

%% ============================================================
%% Section formatting
%% ============================================================
\titleformat{\section}
  {\large\bfseries\color{seriesblue}}
  {}
  {0pt}
  {}
  [\vspace{-2pt}{\color{seriesblue}\rule{\linewidth}{0.7pt}}\vspace{2pt}]

\titleformat{\subsection}
  {\normalsize\bfseries\color{sectiongray}}
  {}
  {0pt}
  {}

\setcounter{secnumdepth}{0}

\setlist[itemize]{
  itemsep=4pt,
  topsep=5pt,
  parsep=0pt,
  label=$\bullet$
}

%% ============================================================
\begin{document}
%% ============================================================

\thispagestyle{empty}

\vspace*{1.2cm}

{\centering
{\color{seriesblue}\fontsize{24}{30}\selectfont\bfseries قسمت دهم}

\vspace{0.5cm}
{\fontsize{15}{20}\selectfont\bfseries سازگاری متقاطع: حل تناقض‌های ظاهری}

\vspace{0.9cm}
{\color{seriesblue}\rule{0.55\linewidth}{1pt}}
\vspace{0.7cm}

{\normalsize
\begin{tabular}{rl}
  \textbf{مجموعه:}     & گزارش پيشرفت مشاور \dash{} بررسی عمیق \\[5pt]
  \textbf{مدت:}        & ۷ تا ۸ دقیقه \\[5pt]
  \textbf{راوی:}       & مجری تنها \\[5pt]
  \textbf{منابع:}      & بخش‌های ۳.۵ و ۵.۲ تا ۵.۴ و ۶ گزارش مشاور \\
\end{tabular}
}

\vspace{0.9cm}
{\color{seriesblue}\rule{0.55\linewidth}{0.4pt}}
\par}

\vspace{1.2cm}

%% ============================================================
\section{افتتاحیه: وقتی اعداد به نظر تعارض دارند}

يک مشاور که گزارش را با دقت می‌خواند چيزهايی پيدا می‌کند که به نظر می‌رسد با هم تضاد دارند. همان کنترل‌گر مقادير لرزش متفاوتی در جداول مختلف دارد. همان اپسيلون با دو مقدار مختلف در دو بخش مختلف ظاهر می‌شود. کنترل‌گر هيبريد اعداد عملکرد نشان می‌دهد با اينکه همه ۱۰۰ اجرای مونت‌کارلو را شکست داد.

اينها خطا نيستند. آن‌ها تفاوت‌های مشروع در زمينه، تعريف معيار، و راه‌اندازی آزمايش هستند. اما بدون درک آشتی، مثل ناسازگاری‌ها به نظر می‌رسند. اين قسمت هر کدام را به صورت سيستماتيک مرور می‌کند.

%% ============================================================
\section{تناقض الف: دو عدد لرزش برای يک کنترل‌گر}

در جدول عملکرد اصلی \lr{MT-5}، کنترل کلاسيک مقدار لرزش ۰.۶۴۷ به علاوه و منهای ۰.۳۴۷ نشان می‌دهد. در جدول آناليز لرزش \lr{FFT}، کنترل کلاسيک شاخص لرزش ۸.۲ نشان می‌دهد. اينها يک معيار نيستند.

مقدار لرزش \lr{MT-5} يک اندازه‌گيری حوزه زمان است: دامنه \lr{RMS} نوسان سيگنال کنترل، محاسبه‌شده در پنجره شبيه‌سازی، ميانگين‌گرفته‌شده در ۱۰۰ اجرا. واحدهايش تقريباً نيوتن است.

شاخص لرزش \lr{FFT} يک اندازه‌گيری حوزه فرکانسی است: نسبت انرژی در باند فرکانس بالا بالای ۲۰ هرتز به کل انرژی سيگنال، بيان‌شده به صورت درصد و مقياس‌شده با ۱۰۰ برای خوانايی.

اينها چيزهای اساساً متفاوتی اندازه می‌گيرند. \lr{RMS} حوزه زمان به شما می‌گويد نوسانات کنترل چقدر بزرگند. شاخص حوزه فرکانسی به شما می‌گويد نوسانات چقدر در محدوده فرکانس بالا متمرکز هستند.

برای ادعای سايش محرک \dash{} ادعا که \lr{STA} سايش سخت‌افزاری را کاهش می‌دهد \dash{} شاخص \lr{FFT} معيار مناسب است. سايش محرک عمدتاً توسط سيکل‌های سوئيچينگ فرکانس بالا هدايت می‌شود، نه توسط دامنه کل حرکت.

%% ============================================================
\section{تناقض ب: نقض بهره \lr{STA} استحکام}

در قسمت‌های ۳ و ۴ علامت‌گذاری کرديم که رديف استحکام \lr{STA} \lr{K-one} مساوی ۲.۰۲ و \lr{K-two} مساوی ۶.۶۷ نشان می‌دهد. اين نقض \lr{K-one} بزرگ‌تر از \lr{K-two} شرط مورنو-اوسوريو مورد نياز برای اثبات همگرايی زمان محدود است.

آيا اين قابل انتشار است؟ پاسخ مستقيم: بدون رفع نه.

بهره‌های اسمی \lr{STA} \dash{} \lr{K-one} مساوی ۸، \lr{K-two} مساوی ۴ \dash{} شرايط نظری را کاملاً ارضا می‌کنند و نتايج معيار مربوطه کاملاً معتبرند. بهره‌های استحکام به عنوان نتايج تجربی که لرزش را در مجموعه ۱۵ سناريويی کاهش می‌دهند برچسب‌گذاری شده‌اند اما تضمين نظری ندارند.

مسير رفع دو گزينه دارد: اجرای مجدد \lr{PSO} استحکام با قيد \lr{K-one} بزرگ‌تر از \lr{K-two} فعال‌شده به صورت محدود در طول جستجو، يا بررسی اينکه آيا يک نتيجه پايداری متفاوت \dash{} شايد يک نتيجه مورنو-اوسوريو گسترده‌شده \dash{} رژيم \lr{K-one} کمتر از \lr{K-two} را پوشش می‌دهد.

%% ============================================================
\section{تناقض ج: اپسيلون ۰.۳ در برابر ۰.۰۲}

در فرمول‌بندی‌های کنترل‌گر و همه معيارهای مونت‌کارلو، اپسيلون مساوی ۰.۳ است. در مطالعه لايه مرزی \lr{MT-6} از قسمت ۸، مقدار نزديک-بهينه ۰.۰۲ يافت شد. کدام درست است؟

هر دو در زمينه‌شان درست هستند.

اپسيلون مساوی ۰.۳ مقدار عملياتی استفاده‌شده در همه نتايج منتشرشده است. بهره‌های \lr{PSO} با اين اپسيلون تنظيم شدند. آمار مونت‌کارلو، اعداد انرژی، زمان‌های يکنواختی، تصديق لياپانوف \dash{} همه از ۰.۳ استفاده کردند.

اپسيلون مساوی ۰.۰۲ يافته مطالعه \lr{MT-6} است: در سه مقدار آزمايش‌شده، اپسيلون ۰.۰۲ شاخص لرزش \lr{FFT} را با دقت رديابی قابل قبول به حداقل رساند. اين نتيجه پس از تکميل معيارها يافت شد.

ناسازگاری زمانی است، نه منطقی. مرحله بعدی صحيح \dash{} مستند در گزارش \dash{} تنظيم مجدد بهره‌های \lr{PSO} با اپسيلون ۰.۰۲ و اجرای مجدد معيارهاست. تا آن زمان، همه نتايج از اپسيلون ۰.۳ استفاده می‌کنند.

%% ============================================================
\section{تناقض د: هيبريد با اعداد فاصله اطمينان با وجود شکست}

گزارش يک رديف برای هيبريد تطبيقی ابر-پيچشی در جدول فاصله اطمينان دارد: فاصله زمان يکنواختی ۱.۷۹ تا ۲.۱۱ ثانيه، فاصله فراجهش ۳.۰ تا ۴.۰ درصد. اما جداول قبلی نشان می‌دهند هيبريد در ۱۰۰ اجرای مونت‌کارلو مقادير شاخص برگرداند، يعنی شکست خورد.

اين اعداد از زمينه‌های آزمايشی متفاوت می‌آيند.

مقادير شاخص \dash{} انرژی ده به توان ششم، لرزش صفر \dash{} به طور خاص از مطالعه \lr{MT-5} با شرايط اوليه مونت‌کارلو استاندارد از قسمت ۶ می‌آيند. اين شرايط اوليه کنترل‌گر هيبريد را به حالت شکست واگرايی بهره فعال می‌کنند.

فاصله‌های اطمينان برای زمان يکنواختی و فراجهش از يک مطالعه اوليه جداگانه می‌آيند \dash{} \lr{QW-2} يا يک اجرای تأييد خاص کنترل‌گر \dash{} جايی که هيبريد با شرايط اوليه کوچک‌تر و کنترل‌شده‌تری آزمايش شد که ناپايداری تقويتی را فعال نکرد. تحت آن شرايط، هيبريد با زمان يکنواختی حدود ۱.۹۵ ثانيه خوب عمل کرد.

هر دو نتيجه با توضيح قسمت ۳ سازگار هستند: هيبريد برای اغتشاشات کوچک کار می‌کند اما برای اغتشاشات بزرگ‌تر شکست می‌خورد. شرايط اوليه \lr{MT-5} به اندازه کافی بزرگ برای فعال کردن شکست بودند. مطالعه پيشين از موارد کوچک‌تری استفاده کرد که نکردند.

%% ============================================================

\begin{mdframed}[
  backgroundcolor=audiotan,
  linecolor=audionote!70,
  linewidth=1pt,
  innerleftmargin=12pt,
  innerrightmargin=12pt,
  innertopmargin=8pt,
  innerbottommargin=8pt
]
\textbf{\color{audionote}[يادداشت صوتی:]}
هيچ‌کدام از تناقض‌های ظاهری در گزارش خطا نيستند. هر کدام تفاوتی در تعريف معيار، زمينه آزمايشی، يا زمان‌بندی آناليز هستند. آشتی هر يک: لرزش دو معيار، يکی حوزه زمان \lr{RMS} و ديگری فرکانسی \lr{FFT}. بهره‌های استحکام \lr{STA}: نتايج تجربی، نه تضمين‌شده نظری. دو مقدار اپسيلون: مقدار عملياتی معيار در برابر يافته مطالعه. اعداد هيبريد: زمينه‌های آزمايشی متفاوت جايی که شرايط اوليه کوچک‌تر ناپايداری را فعال نکردند.
\end{mdframed}

\vspace{0.6cm}

%% ============================================================
\section{نتيجه‌گيری}

هر چهار تناقض ظاهری به يک مسئله زمينه‌ای حل می‌شوند: يک کلمه يا عدد يکسان ممکن است بسته به زمينه اندازه‌گيری معانی متفاوتی داشته باشد. وظيفه گزارش اين تفاوت‌های زمينه‌ای را صريح کردن به جای اجازه دادن به ظاهر شدن به عنوان تناقض است. و اين کار را انجام می‌دهد \dash{} هر مورد يک يادداشت، يک برچسب، يا يک بخش تواضع‌دهنده روشن‌سازی تفاوت دارد.

قسمت بعدی و يازدهم: آنچه هنوز بايد نوشته شود قبل از ارسال نهايی پايان‌نامه.

%% ============================================================
\vspace{1.5cm}
\begin{center}
{\color{seriesblue}\rule{0.45\linewidth}{0.4pt}}\\[8pt]
{\small\color{sectiongray}
منابع گزارش: بخش‌های ۳.۵ و ۵.۲ تا ۵.۴ و ۶
}
\end{center}

\end{document}
